\section{Does field-level soil-temperature sensing change the schedule?}
\label{sec:iot}

\subsection{The gap the sensor would fill}

CERES-Maize does not run the phase that sets flowering date on air
temperature. While \texttt{LEAFNO .LE. 10}, \texttt{MZ\_PHENOL.for} (lines
457--496) forms daily thermal time from \emph{soil} temperature, estimated
from air temperature and solar radiation:
\begin{align}
a &= 0.01061\,\mathrm{SRAD} + 0.5902, \label{eq:acoef}\\
T_{\mathrm{day}} &= a\,T_{\max} + (1-a)\,T_{\min}, \label{eq:tdsoil}\\
T_{\mathrm{night}} &= 0.36354\,T_{\max} + 0.63646\,T_{\min}, \label{eq:tnsoil}\\
\Delta T\!T &= \tfrac{1}{2}\!\left[\max(T_{\mathrm{night}},T_b)
  + \min(T_{\mathrm{day}},D_{\mathrm{opt}})\right] - T_b, \label{eq:dtt}
\end{align}
with $T_b=8.0$ and $D_{\mathrm{opt}}=34.0$\,\textdegree C from the shipped
ecotype file, $\Delta TT$ then capped at $D_{\mathrm{opt}}-T_b$. The clamps
in~\eqref{eq:dtt} apply to each component \emph{before} the average: a
$40/20$\,\textdegree C pair yields $19.0$, averaging first $22.0$.
Crucially, \eqref{eq:acoef}--\eqref{eq:tnsoil} contain no field-specific
term, so two fields under one weather station receive identical estimated
soil temperature whatever their soil does. That imputation is the gap a
buried probe would close, and the only point where a soil-temperature
measurement has anywhere to enter.

\subsection{Simulation design}

\textbf{Nothing here is a measurement.} No soil-temperature observation was
obtained for this study. Paired 5\,cm bare and vegetated probes exist in the
Oklahoma Mesonet but sit behind an interactive request form and a
robots-disallowed endpoint; USCRN publishes no probe-level soil temperature,
its readme files stating the reported values are ``calculated from these
multiple independent measurements''. The sensing layer is explicitly
simulated and every parameter of it is recorded on each result row.

We implement \eqref{eq:acoef}--\eqref{eq:dtt} directly and calibrate one
thermal threshold per cultivar against the existing DSSAT table ($n=1281$
per cultivar), reproducing it to $2.62$\,d RMSE with mean absolute bias
$0.14$\,d against $2.70$\,d for the same arithmetic on air temperature.
This is a self-consistent substrate, not a claim to reproduce DSSAT: truth
and all planners run on it, so its absolute error cancels in the paired
comparison and its bias --- which would not --- is near zero.

Each field receives a constant soil thermal offset $\delta_i$. The default
scale $\sigma=0.7$\,\textdegree C is \emph{not} chosen here: it is the value
already carried by \texttt{sim.core.make\_farm}, committed before this
analysis existed. $\sigma$ is swept from $0$ regardless, and $\sigma=0$ is a
null scenario in which a sensor must show no gain. A virtual probe reads the
field's soil temperature with Gaussian noise of $0.5$\,\textdegree C --- the
documented accuracy of the Mesonet thermistors whose data we could not
obtain --- and misses days at a configurable rate; $\hat\delta_i$ is the
mean over observed days of the residual against what
\eqref{eq:tdsoil}--\eqref{eq:tnsoil} would impute. Observation closes on day
273, before the window opens on day 274, so the channel is
non-anticipative: a field offset is a persistent property, not a forecast.

Three planners differ only in belief: \textsc{conventional} believes every
offset is zero, exactly what CERES does; \textsc{sensor} believes
$\hat\delta_i$, falling back to zero where estimation failed;
\textsc{oracle} believes the truth and bounds the gain. Everything else ---
fields, parent pairs, areas, weather years, capacity model and budget,
response scenario, objective, candidate enumeration --- is the same object
for all three, and every coefficient and executed score passes through one
scoring function.

\subsection{Capacity must bind before the question is meaningful}

At the farm size used elsewhere here the capacity constraint is almost
inactive --- at $N=12$ the constrained optimum is within $0.06\%$ of the
unconstrained one --- so measuring the value of information there measures
nothing about scarcity. At $N=60$ it binds properly, $17.7\%$ of objective
under the visits model, with only $53$ of $60$ fields sown. All
capacity-constrained results use $N=60$.

\subsection{Results}

\begin{table}[t]
\centering
\caption{Change in objective against the conventional CERES pathway.
Capacity-constrained rows: visits model, $N=60$, 6 farm seeds per level
($\sigma=0.7$ repeated at 12 seeds gives $+0.583\%$ $[+0.435,+0.732]$).
95\% confidence intervals across seeds.}
\label{tab:iot}
\begin{tabular}{lrrrrr}
\toprule
$\sigma$ (\textdegree C) & sensor (\%) & 95\% CI & seeds $>0$ & oracle (\%) & captured \\
\midrule
\multicolumn{6}{l}{\emph{capacity-constrained, $N=60$}}\\
0.00 & $-1.410$ & $[-1.505,-1.316]$ & 0/6 & $+0.000$ & --- \\
0.50 & $+0.455$ & $[+0.296,+0.615]$ & 6/6 & $+2.005$ & 0.23 \\
0.70 & $+0.563$ & $[+0.392,+0.735]$ & 6/6 & $+2.279$ & 0.25 \\
1.00 & $+1.050$ & $[+0.836,+1.263]$ & 6/6 & $+2.792$ & 0.38 \\
2.00 & $+2.125$ & $[+1.789,+2.462]$ & 6/6 & $+3.891$ & 0.55 \\
4.00 & $+6.065$ & $[+2.972,+9.157]$ & 6/6 & $+7.493$ & 0.81 \\
\bottomrule
\end{tabular}
\end{table}

\textbf{Scarcity makes the information worth having.} Matched at $N=60$ and
$\sigma=0.7$, the sensor gains $+0.191\%$ $[-0.008,+0.389]$ without a
binding capacity constraint --- an interval spanning zero --- and
$+0.583\%$ $[+0.435,+0.732]$ with one, positive in all twelve seeds; the
oracle rises from $+1.642\%$ to $+2.279\%$ over the same contrast.
Information about which fields run warm is worth more when the crew cannot
serve every field, because it then also informs \emph{which} fields to sow.

\textbf{Scarcity also makes noise more expensive.} At $\sigma=0$ the oracle
gain is exactly zero, as it must be, while the sensor loses $1.410\%$
$[-1.505,-1.316]$, negative in every seed, against $0.164\%$ unconstrained.
Noise read as real offset does not merely mis-set an offset; under scarcity
it misallocates a day the crew cannot get back.

\textbf{A realistic probe leaves three quarters of the value unclaimed.} A
noise ablation at $\sigma = 0.7$, six seeds, isolates why: a noiseless probe
recovers the oracle \emph{exactly} ($+2.279\%$ against $+2.279\%$), and the
captured share falls to $0.41$ at $0.25$\,\textdegree C, $0.25$ at
$0.5$\,\textdegree C, $0.20$ at $1.0$ and $0.15$ at $2.0$. The limiting
quantity is instrument accuracy, not the scheduler and not the observation
campaign. Estimation degrades gracefully by comparison: losing three
quarters of the observations moves the captured share only from $0.25$ to
$0.18$, and with every observation lost the gain is exactly $+0.000\%$ with
no plan changed, confirming the zero-belief fallback.



\textbf{Solver validity.} Every capacity-constrained cell above was solved
on the visits model with measured LP gap median $0.0001$, maximum $0.0003$,
inside the $0.005$ acceptance rule, and no arm hit the time limit. The
operations, hours and hectares models do not reach that standard at $N=60$
and more solver time does not help: at a 45\,s limit every arm terminated
time-limited with gaps $0.011$--$0.034$, and raising the limit to $300$\,s
left the gaps on the same cell unchanged to five decimals
($0.01140$, $0.02253$, $0.01441$) at twenty times the cost. The limitation
is LP bound quality for these cost structures, not search time. Those cells
are lower bounds, their spread exceeds the effect measured, and one produced
a negative oracle gain that cannot be real, so they are excluded from
inference. Robustness to capacity accounting is therefore \emph{not}
established here.

\subsection{What this does and does not establish}

It establishes, inside this model and with simulated sensing, that the
decision value of per-field soil-temperature measurement is governed by the
ratio of between-field heterogeneity to probe accuracy and by how tightly
sowing capacity binds. Both raise the gain, and scarcity raises the cost of
error at the same time. That is a transferable criterion for when such a
system is worth building.

It does not establish what the between-field heterogeneity of a commercial
seed-production block actually is. Published same-site contrasts at 5\,cm
run 1--6\,\textdegree C between deliberately imposed treatments, but such a
block is managed uniformly, and the one continuous relationship available
--- soil temperature against percent residue cover --- implies
$0.11$--$0.59$\,\textdegree C between fields under one protocol. At the top
of that range the constrained sensor gains about $+0.46\%$; at the bottom it
is indistinguishable from zero and, if the fields are truly identical,
actively harmful. Establishing the real figure requires paired measurements
we were unable to obtain. The emulator's 2.6\,d RMSE against DSSAT also
exceeds the one-day decision unit; the paired comparison cancels this, but
it would not cancel for any absolute claim about flowering date. \textbf{No seed-yield
improvement is claimed and no physical sensor was deployed or tested}: the
objective is this study's own area-weighted relative-yield response to the
anthesis--silking interval.
